\chapter{Frais, remises et économie des plateformes}\label{ch:mx:fees-rebates-and-venue-economics}

Le même titre se négocie au même prix sur deux bourses. L'une verse une remise à l'ordre au repos et compte vingt mille actions en file devant lui~; l'autre
facture des frais à l'ordre au repos et n'en compte que huit cents. Un trader qui veut être exécuté doit payer pour se placer dans la file la plus courte. Les frais de bourse se chiffrent à quelques
centièmes de centime par action, mais la fourchette d'une valeur liquide vaut un centime, et la répartition des frais entre le côté qui se met au repos et celui qui prend
modifie où vont les ordres, la longueur des files et qui se fait sélectionner. Ce chapitre évalue le prix des ordres nets de frais, montre quand la répartition d'un frais compte,
et fait tourner côte à côte deux plateformes du simulateur de bourse, l'une maker-taker et l'autre inversée, pour mesurer ce qu'une remise rapporte.

\section{Maker-taker, inversé et tarification uniforme}

Sous une tarification maker-taker (livre~1, chapitre~9), la plateforme verse au côté qui se met au repos une remise et facture au côté preneur des frais~; une plateforme inversée
(taker-maker) fait l'inverse. Les deux encaissent la différence.

\begin{definition}[Tarification uniforme]\label{def:mx:fees-rebates-and-venue-economics:flat}
Sous une \emph{tarification uniforme}\index{tarification uniforme} (flat pricing), une plateforme facture les deux côtés d'une transaction, ou leur facture le même montant par action quel que soit leur rôle,
de sorte que se mettre au repos et prendre la liquidité coûtent la même chose.
\end{definition}

La \omterm{def:mx:fees-rebates-and-venue-economics:flat}{tarification uniforme} est la norme pour les contrats à terme, où la bourse facture à chaque côté des frais par contrat, et sur de nombreux marchés d'actions hors des
États-Unis~; la tarification maker-taker et la tarification inversée sont la façon dont les bourses d'actions américaines se disputent les deux côtés du flux d'ordres, dans la limite d'un plafond des frais d'accès aux
cotations affichées. Les deux plateformes Cboe de l'encadré servent d'exemple tout au long du chapitre~; la simulation utilise leurs tarifs standard.

\begin{dated}{2026-09}{Deux barèmes de frais et le nouveau plafond}\label{dat:mx:fees-rebates-and-venue-economics:cboe}
Tarifs standard pour les titres à 1.00~\$ ou plus, en vigueur le 1\up{er} septembre 2026. Cboe BZX (maker-taker)~: une remise de 0.0016~\$ par action pour l'ajout
de liquidité affichée, des frais de 0.0030~\$ pour le retrait. Cboe BYX (taker-maker)~: des frais de 0.0020~\$ pour l'ajout, une remise de 0.0002~\$ pour le retrait.
Les deux barèmes prévoient des paliers et des codes de frais qui modifient ces tarifs pour de nombreux membres. Les modifications de 2024 apportées à la Rule 610(c) par la SEC abaissent le plafond des frais d'accès
à une cotation protégée affichée à 1.00~\$ ou plus, de 0.0030~\$ à 0.001~\$ par action~; des ordonnances d'exemption des 31 octobre 2025 et 11 juin 2026 ont reporté la mise en conformité au
premier jour ouvré de novembre 2027.
\end{dated}

\section{Prix nets des frais et pas de cotation effectif}

\begin{definition}[Prix net des frais]\label{def:mx:fees-rebates-and-venue-economics:adjusted}
Le \emph{prix net des frais}\index{prix net des frais} (fee-adjusted price) d'une transaction est son prix majoré des frais par action pour un acheteur et minoré de ces frais pour un vendeur, une remise
comptant comme des frais négatifs~: ce que la transaction coûte ou rapporte une fois la plateforme payée. Avec $f$ les frais liés au rôle dans la transaction et $\varepsilon=\pm1$ son côté,
il vaut $p+\varepsilon f$.
\end{definition}

Un acheteur qui prend 1\,000 actions à 20.00 paie 20.0030~\$ par action sur la plateforme maker-taker et 19.9998~\$ sur la plateforme inversée~: 3.20~\$ de moins
sur la plateforme inversée, au même cours affiché. Les différences de frais créent des prix entre les ticks.

\begin{definition}[Pas de cotation effectif]\label{def:mx:fees-rebates-and-venue-economics:efftick}
Le \emph{pas de cotation effectif}\index{pas de cotation effectif} (effective tick size) est la plus petite différence entre les \omterm{def:mx:fees-rebates-and-venue-economics:adjusted}{prix nets des frais} auxquels un côté peut traiter
lorsque le même instrument est coté sur des plateformes aux frais différents~: avec des frais preneurs $f_1$ et $f_2$ et un tick $\delta_{\mathrm{tick}}$, les offres
nettes $p+f_1$ et $p+f_2$ s'entrelacent avec des écarts $|f_1-f_2|$ et $\delta_{\mathrm{tick}}-|f_1-f_2|$ (frais réduits modulo le tick).
\end{definition}

Avec les deux frais preneurs de Cboe, 0.30 et $-0.02$ centime, les offres nettes d'une valeur à un centime se trouvent alternativement à 0.32 et 0.68 centime d'écart~: un cours affiché sur la
plateforme inversée constitue une amélioration de prix de 0.32 centime que le tick d'un centime interdit sur une plateforme isolée (voir la \cref{fig:mx:fees-rebates-and-venue-economics:ladder}).
Chao, Yao et Ye (2019) soutiennent que c'est pour cette raison qu'un tick discret fragmente les bourses d'actions américaines et disperse leurs structures de frais~: les plateformes
se différencient par les frais parce qu'elles ne peuvent pas le faire par le prix.

\begin{omfigure}
\begin{tikzpicture}[font=\footnotesize,x=6.5cm]
\draw[->,gray] (-0.12,0) -- (1.5,0);
\node[gray,below] at (0.75,-0.9) {offre nette, en centimes au-dessus de 20.00};
\foreach \x in {0,0.5,1}{\draw[gray] (\x,0.08) -- (\x,-0.08) node[below] {\x};}
\foreach \x/\l in {-0.02/INV à 20.00,0.98/INV à 20.01}{\fill[omThm] (\x,0) circle (2.5pt); \node[above,omThm] at (\x,0.12) {\l};}
\foreach \x/\l in {0.30/MT à 20.00,1.30/MT à 20.01}{\fill[omDef] (\x,0) circle (2.5pt); \node[above,omDef] at (\x,0.12) {\l};}
\draw[<->] (-0.02,0.75) -- (0.30,0.75) node[midway,above] {0.32};
\draw[<->] (0.30,0.75) -- (0.98,0.75) node[midway,above] {0.68};
\draw[<->] (0.98,0.75) -- (1.30,0.75) node[midway,above] {0.32};
\end{tikzpicture}

\omcaption{Offres nettes d'une valeur cotée à 20.00 et 20.01 sur une plateforme maker-taker (frais preneur de 0.30 centime) et une plateforme inversée (remise preneur
de 0.02 centime), en centimes au-dessus de 20.00. La grille vue par un preneur est plus fine que le tick d'un centime.}\label{fig:mx:fees-rebates-and-venue-economics:ladder}
\end{omfigure}

Un preneur compare les plateformes sur les \omterm{def:mx:fees-rebates-and-venue-economics:adjusted}{prix nets des frais}, et un routeur répartit un ordre marketable entre elles dans cet ordre.

\omcode{code/firm/venuefees/firm_venuefees.py}{40}{51}{Un ordre marketable réparti sur les cotations affichées des plateformes, par prix net des frais croissant, le meilleur en premier.}\label{lst:mx:fees-rebates-and-venue-economics:route}

\section{La répartition des frais est-elle neutre~?}

\begin{definition}[Neutralité des frais]\label{def:mx:fees-rebates-and-venue-economics:neutral}
La répartition des frais d'une plateforme entre teneurs de marché et preneurs est \emph{neutre}\index{neutralité des frais} (neutral) lorsqu'on déplace une partie des frais d'un côté à l'autre,
le total $F=f^{make}+f^{take}$ étant fixé, laisse inchangés les \omterm{def:mx:fees-rebates-and-venue-economics:adjusted}{prix nets des frais} de chaque trader et les résultats du marché~: ce sont les cotations qui absorbent le déplacement.
\end{definition}

\begin{proposition}[Neutralité sans tick]\label{prop:mx:fees-rebates-and-venue-economics:neutral}
Supposons que des fournisseurs de liquidité en concurrence affichent l'offre $a$ à laquelle une vente est à l'équilibre, soit $a-f^{make}=R$, où $R$ est leur prix de réserve (valeur, coûts
et sélection adverse). Le prix net du preneur vaut $a+f^{take}=R+F$, et le prix net du teneur vaut $R$~: les deux ne dépendent que des frais totaux. Avec un
tick $\delta_{\mathrm{tick}}$, l'offre doit valoir $\lceil(R+f^{make})/\delta_{\mathrm{tick}}\rceil\,\delta_{\mathrm{tick}}$, et tout changement de répartition
qui n'est pas un multiple du tick modifie les \omterm{def:mx:fees-rebates-and-venue-economics:adjusted}{prix nets des frais}.
\end{proposition}

\begin{proof}
Sans tick, $a=R+f^{make}$ et $a+f^{take}=R+F$~; déplacer $\Delta$ de $f^{take}$ vers $f^{make}$ déplace $a$ de $+\Delta$ et rien d'autre. Avec un
tick, $a$ est le plus petit prix de la grille supérieur ou égal à $R+f^{make}$, si bien que le teneur conserve une rente $a-f^{make}-R\in[0,\delta_{\mathrm{tick}})$ qui dépend
de $f^{make}$ modulo le tick.
\end{proof}

Déplacer 0.20 centime des frais preneurs (de 0.30 à 0.10) vers le côté teneur (une remise de 0.16 devient des frais de 0.04) montre la différence. Sans tick, une offre à
10.01 passe à 10.012~: le preneur paie toujours 10.013 tout compris et le teneur touche toujours 10.0116 net. Sur une grille à un centime, l'offre doit rester à 10.01 (le preneur paie désormais
10.011, le teneur perd 0.2 centime) ou saute à 10.02 (le preneur paie 10.021). Lorsque les cotations ne peuvent pas bouger, la rente passe à la position dans la file~:
les teneurs se disputent en rejoignant les files, et la répartition des frais détermine leur longueur. Colliard et Foucault (2012) modélisent comment le niveau et la répartition des frais
modifient les taux d'exécution et les fourchettes~; Foucault, Kadan et Kandel (2013) expliquent la tarification maker-taker par la vitesse à laquelle les deux côtés réagissent. À la Bourse de Toronto
Malinova et Park (2015) constatent que les cotations affichées ne s'ajustent pas à une modification de la répartition des frais ni aux coûts des demandeurs de liquidité, frais inclus,
tandis que les \omterm{def:mx:decomposing-the-spread:quoted}{fourchettes affichées} se resserrent et que les ordres agressifs deviennent plus fréquents.

\section{Paliers et activité de la plateforme}

Le revenu d'une plateforme tiré du trading est la différence entre ce que paient les preneurs et ce que reçoivent les teneurs, multipliée par son volume~; les tarifs de ses membres dépendent des
paliers de volume (livre~1, chapitre~29), qui accordent les plus grandes remises aux plus gros fournisseurs de liquidité et font des seuils de volume des falaises.
Les frais de négociation ne constituent qu'une part de l'activité~: les données de marché et la connectivité sont vendues aux mêmes membres.

La simulation construit le choix que créent les paliers et les frais. Deux moteurs de \texttt{firm.exchsim} font tourner une même valeur avec un tick d'un centime~: MT avec les tarifs standard de BZX,
INV avec ceux de BYX. Un petit marché réactif fournit le flux~: un prix efficient $P^\ast$ qui se déplace d'un tick à des instants aléatoires~; des \omterm{def:mx:the-limit-order-book:orders}{ordres à cours limité}
qui rejoignent le meilleur cours ou se placent à quelques ticks derrière et s'annulent à taux constant, et aussitôt avec probabilité un demi lorsque $P^\ast$ franchit
leur prix~; des \omterm{def:mx:the-limit-order-book:orders}{ordres au marché} de bruit~; et des \omterm{def:mx:the-limit-order-book:orders}{ordres au marché} informés à un taux proportionnel à la distance entre $P^\ast$ et le milieu. Chaque \omterm{def:mx:the-limit-order-book:orders}{ordre à cours limité} va sur MT
ou sur INV avec probabilité un demi. Un preneur sensible aux frais route selon le \omterm{def:mx:fees-rebates-and-venue-economics:adjusted}{prix net des frais} (voir le \cref{lst:mx:fees-rebates-and-venue-economics:route}), si bien qu'à prix égaux
il prend d'abord la file de INV~; un preneur insensible aux frais prend d'abord la file affichée la plus grande, comme le ferait un routeur qui ne regarde que la taille. Chaque cas tourne
sur trois graines de deux heures.

\omcode{code/microstructure/07-fees-rebates-and-venue-economics/python/mx_fees.py}{128}{137}{Les preneurs du marché à deux plateformes~: routage fee-aware par prix net des frais, routage fee-blind par taille affichée.}\label{lst:mx:fees-rebates-and-venue-economics:takers}

Lorsque tous les preneurs sont fee-aware, la plateforme inversée est la tête d'une file consolidée et la plateforme maker-taker en est la queue. INV capte 77~\% du
volume. Ses \omterm{def:mx:the-limit-order-book:orders}{ordres à cours limité} reçoivent une exécution pour 28.2~\% des actions affichées, après un délai médian de 4.6 secondes, contre 8.4~\% après 8.8 secondes sur MT, où la
file au meilleur cours est plus longue (910 actions par côté en moyenne, contre 510). Les exécutions sur MT interviennent lorsque la file d'INV est épuisée, c'est-à-dire lorsque le flux est
important ou informé~: le milieu se déplace de 0.63 cents en cas d'exécution sur MT dans les 30 secondes, contre 0.13 en cas d'exécution sur INV (erreurs types de 0.04). La capture est la même
sur les deux, soit 0.54 cents. Net de la sélection adverse, une exécution sur MT vaut $-0.09$ cents avant la rémunération et une exécution sur INV vaut 0.40 cents avant les frais.

\begin{omfigure}
\begin{tikzpicture}
\begin{groupplot}[group style={group size=2 by 1,horizontal sep=1.5cm},width=6.6cm,height=5cm,
  tick label style={font=\footnotesize},label style={font=\small},grid=major,grid style={gray!25},table/col sep=comma,
  xtick={0,0.5,1},xlabel={part des preneurs fee-aware},xmin=-0.1,xmax=1.1]
\nextgroupplot[title={\small part exécutée des actions affichées},ymin=0,ymax=0.32,
  legend columns=2,legend style={font=\footnotesize,at={(1.1,-0.3)},anchor=north,draw=none,fill=none,
  /tikz/every even column/.append style={column sep=8pt}}]
\addplot[omDef,very thick,mark=*] table[x=aware,y=fill_mt]{figdata/microstructure/07-fees-rebates-and-venue-economics/venues.csv};
\addplot[omThm,very thick,mark=square*] table[x=aware,y=fill_inv]{figdata/microstructure/07-fees-rebates-and-venue-economics/venues.csv};
\legend{maker-taker (MT),inversée (INV)}
\nextgroupplot[title={\small sélection adverse à 30 secondes (cents)},ymin=0,ymax=0.75]
\addplot[omDef,very thick,mark=*] table[x=aware,y=adverse_mt]{figdata/microstructure/07-fees-rebates-and-venue-economics/venues.csv};
\addplot[omThm,very thick,mark=square*] table[x=aware,y=adverse_inv]{figdata/microstructure/07-fees-rebates-and-venue-economics/venues.csv};
\end{groupplot}
\end{tikzpicture}

\omcaption{Deux plateformes, même flux, la moitié des \omterm{def:mx:the-limit-order-book:orders}{ordres à cours limité} sur chacune. Plus les preneurs routent par \omterm{def:mx:fees-rebates-and-venue-economics:adjusted}{prix net des frais}, plus la plateforme inversée est exécutée
en premier et plus les exécutions de la plateforme maker-taker proviennent d'un flux qui continue de déplacer le prix. Moyennes de trois graines de deux heures. Données~:
\texttt{mx\_fees.by\_awareness}.}\label{fig:mx:fees-rebates-and-venue-economics:venues}
\end{omfigure}

Que rapporte une action affichée~? Sa valeur est la part exécutée multipliée par la capture, moins la sélection adverse, moins les frais de création. Lorsque tous les preneurs sont fee-aware, une
action affichée sur MT vaut 0.006 cent et une sur INV 0.057 (l'écart a une erreur type de 0.009)~: la rémunération ne paie pas l'arrière de la
file. Pour qu'un fournisseur soit indifférent, MT devrait verser une rémunération de 0.77 cents, soit plus du double de ses frais de prise de 0.30 cent~: aucune plateforme maker-taker ne peut
acheter cette file. Lorsque la moitié des preneurs sont fee-blind, la file de MT est touchée en premier par ceux qui regardent la taille, et le classement s'inverse~: 0.052 cents sur MT contre
0.028 sur INV (erreur type de 0.006), et les frais de création d'indifférence sur MT sont proches de zéro (0.006 cent), si bien que la rémunération de 0.16 cent dépasse ce dont MT a besoin. Sans aucun
preneur fee-aware, une action affichée sur MT vaut 0.074 contre 0.015. La manière dont les preneurs routent détermine ce qu'achète une rémunération.

\section{Évaluation à partir des expériences sur les frais}

La SEC a tenté de mesurer les frais comme le Tick Size Pilot avait mesuré le tick. Elle a adopté le Transaction Fee Pilot, Rule 610T, le 19 décembre 2018~: 1\,460
actions tirées au hasard dans deux groupes de test, l'un avec un plafond de 0.0010~USD sur les frais de transaction en bourse (le plafond alors en vigueur, fixé en 2005, était de
0.0030~USD) et l'autre dans lequel les bourses ne pouvaient pas verser de rémunérations, ainsi que toutes les autres actions d'un groupe témoin. Les bourses ont contesté la règle, et la D.C. Circuit
l'a annulée le 16 juin 2020, avant sa mise en œuvre. Les données probantes proviennent donc de modifications que les plateformes ont apportées d'elles-mêmes et de données de routage. L'étude de Malinova
et Park sur Toronto en est une~; Battalio, Corwin et Jennings (2016) ont constaté que des courtiers de détail envoyaient des \omterm{def:mx:the-limit-order-book:orders}{ordres à cours limité} aux plateformes versant les plus fortes rémunérations
et une relation négative entre la qualité d'exécution des \omterm{def:mx:the-limit-order-book:orders}{ordres à cours limité} et le niveau de la rémunération, le conflit que décrivaient Angel, Harris et Spatt. Leur constat est
le versant empirique de la simulation~: un ordre placé pour sa rémunération se retrouve à l'arrière de la file consolidée.

\section{Tutoriel~: la rémunération qui a acheté une file}\label{tut:mx:fees-rebates-and-venue-economics:tut}

\begin{tutorial}
\textbf{Objectif.} Fixer les prix des ordres nets des frais, puis mesurer sur deux plateformes simulées ce qu'achète une rémunération.
\textbf{État final~:} \cref{fig:mx:fees-rebates-and-venue-economics:venues} et les chiffres des sections 2 à 4.
\begin{tutsteps}
\item \textbf{\omterm{def:mx:fees-rebates-and-venue-economics:adjusted}{Prix nets des frais}.} \texttt{firm\_venuefees.fee\_adjusted} et \texttt{effective\_tick} avec les taux de l'encadré.
\item \textbf{Neutralité.} \texttt{neutral\_ask(ask, take, take\_new, tick)} avec et sans \omterm{def:mx:the-limit-order-book:tick}{pas de cotation}.
\item \textbf{Deux plateformes.} \texttt{mx\_fees.two\_venues(share\_mt, fees, seconds, seed, flow=Flow(aware=\dots))}~: pour chaque plateforme, la part exécutée, le délai jusqu'à la
première exécution, la capture, la sélection adverse, la valeur par action affichée, la file et la part de volume.
\item \textbf{Prise en compte des frais.} \texttt{by\_awareness()} moyenne trois graines pour 0, une moitié et tous les preneurs fee-aware~; \texttt{indifference\_make} donne le
frais de création qui égalisent les deux plateformes~; tracer avec \texttt{fig\_fees.py}.
\end{tutsteps}
\textbf{Pour aller plus loin.} Appliquer le nouveau plafond de 0.10 cent aux frais de prise de MT et baisser sa rémunération pour conserver sa marge~; laisser chaque fournisseur choisir la plateforme
offrant la meilleure valeur espérée compte tenu des files qu'il observe, et regarder les parts se déplacer.
\end{tutorial}

\section{Construction~: frais et choix de plateforme}\label{bld:mx:fees-rebates-and-venue-economics:venuefees}

\begin{build}
\bfield{Rôle} \omterm{def:mx:fees-rebates-and-venue-economics:adjusted}{Prix nets des frais} et choix de plateforme pour le routeur d'ordres intelligent du chapitre~18, l'analyse des coûts de transaction du chapitre~19 et
les teneurs de marché du livre~11.
\bfield{Interface} \texttt{fee\_adjusted(price, side, fee)}, \texttt{effective\_tick(tick, fees)}, \texttt{route\_take(quotes, side, qty)},
\texttt{passive\_value(fill, capture, adverse, make)}, \texttt{indifference\_make(\dots)}, \texttt{neutral\_ask(ask, take, take\_new, tick)},
\bfield{Interface} \texttt{tier\_make\_fee(schedule, adav, tcv)} sur \texttt{firm.feesched} du livre~1.
\bfield{Règles} Frais par action, positifs lorsqu'ils sont payés~; un acheteur ajoute les frais, un vendeur les retranche~; le routage trie les cotations par \omterm{def:mx:fees-rebates-and-venue-economics:adjusted}{prix net des frais}, sur la
seule taille affichée.
\bfield{Tests d'acceptation} \texttt{code/firm/venuefees/tests/}~: \omterm{def:mx:fees-rebates-and-venue-economics:adjusted}{prix nets des frais} et un achat et une vente séparés à la main~; le \omterm{def:mx:fees-rebates-and-venue-economics:efftick}{pas de cotation effectif} pour deux et
trois frais~; les frais de création indifférents résolvent leur équation~; le cours vendeur \omterm{def:mx:fees-rebates-and-venue-economics:neutral}{neutre} avec et sans \omterm{def:mx:the-limit-order-book:tick}{pas de cotation}~; des frais de création à paliers.
\bfield{Pour aller plus loin} Routage avec probabilités d'exécution et positions dans la file (\texttt{firm.queuevalue} du chapitre~6)~; planification mensuelle des paliers avec les seuils de
\texttt{firm.feesched}.
\end{build}

\begin{omsources}
\item J.-E.~Colliard et T.~Foucault, ``Trading fees and efficiency in limit order markets'', \emph{Review of Financial Studies} 25(11), 2012.
\item T.~Foucault, O.~Kadan et E.~Kandel, ``Liquidity cycles and make/take fees in electronic markets'', \emph{Journal of Finance} 68(1), 2013.
\item K.~Malinova et A.~Park, ``Subsidizing liquidity: the impact of make/take fees on market quality'', \emph{Journal of Finance} 70(2), 2015.
\item J.~J.~Angel, L.~E.~Harris et C.~S.~Spatt, ``Equity trading in the 21st century: an update'', \emph{Quarterly Journal of Finance} 5(1), 2015.
\item R.~Battalio, S.~A.~Corwin et R.~Jennings, ``Can brokers have it all? On the relation between make-take fees and limit order execution quality'',
\emph{Journal of Finance} 71(5), 2016.
\item S.~Chao, C.~Yao et M.~Ye, ``Why discrete price fragments U.S. stock exchanges and disperses their fee structures'', \emph{Review of Financial
Studies} 32(3), 2019.
\item United States Court of Appeals for the District of Columbia Circuit, \emph{New York Stock Exchange LLC v. SEC}, No. 19-1042, 2020.
\item Cboe BZX and BYX U.S. Equities fee schedules, septembre 2026.
\end{omsources}

\section{Exercices}

\begin{exercise}[$\star$]\label{exo:mx:fees-rebates-and-venue-economics:1}
Calculer le \omterm{def:mx:fees-rebates-and-venue-economics:adjusted}{prix net des frais} d'un achat de 1\,000 actions à 20.00 avec les frais de prise BZX de l'encadré, puis avec la rémunération de prise de BYX, et déterminer de combien le choix
de la plateforme permet d'économiser.
\end{exercise}

\begin{exercise}[$\star$]\label{exo:mx:fees-rebates-and-venue-economics:2}
Déterminer le \omterm{def:mx:fees-rebates-and-venue-economics:efftick}{pas de cotation effectif} d'une action dont le \omterm{def:mx:the-limit-order-book:tick}{pas de cotation} est d'un cent et négociée sur deux plateformes dont les frais de prise valent 0.30 et $-0.02$ cents~?
\end{exercise}

\begin{exercise}[$\star$]\label{exo:mx:fees-rebates-and-venue-economics:3}
Une action affichée sur INV est exécutée dans 30~\% des cas, capture 0.5 cents et perd 0.1 au titre de la sélection adverse, avec des frais de création de 0.20 cent~; sur MT, elle est exécutée dans 10~\% des
cas, capture 0.5 cents et perd 0.6, avec une rémunération de 0.16 cent. Calculer la valeur de chacune par action affichée~?
\end{exercise}

\begin{exercise}[$\star\star$]\label{exo:mx:fees-rebates-and-venue-economics:4}
Dans l'exercice~3, quels frais de création sur MT rendraient un fournisseur indifférent~? Une plateforme dont les frais de prise sont plafonnés à 0.30 cents peut-elle les verser~?
\end{exercise}

\begin{exercise}[$\star\star$]\label{exo:mx:fees-rebates-and-venue-economics:5}
Une plateforme transfère 0.20 cent de ses frais de prise de 0.30 cent vers le côté création, où une rémunération de 0.16 cent devient des frais de 0.04 cent. Que devient un cours vendeur à 10.01,
son \omterm{def:mx:fees-rebates-and-venue-economics:adjusted}{prix net des frais} pour le preneur de liquidité et son prix net pour le fournisseur de liquidité, sans tick et avec un tick d'un cent~?
\end{exercise}

\begin{exercise}[$\star\star$]\label{exo:mx:fees-rebates-and-venue-economics:6}
Pourquoi la sélection adverse sur la plateforme maker-taker augmente-t-elle avec la part de preneurs sensibles aux frais, de 0.31 à 0.63 cent, alors que la capture reste à 0.54~?
\end{exercise}

\begin{exercise}[$\star\star\star$]\label{exo:mx:fees-rebates-and-venue-economics:7}
\emph{Programmation.} Exécuter \texttt{two\_venues} avec les frais de prise de MT fixés au nouveau plafond de 0.10 cent et sa rémunération ramenée à 0.06 cent (même marge diminuée de 0.10). Avec tous
les preneurs sensibles aux frais, comment évoluent la part exécutée de MT, la sélection adverse et la valeur par action affichée, et pourquoi~?
\end{exercise}

\begin{exercise}[$\star\star\star$]\label{exo:mx:fees-rebates-and-venue-economics:8}
\emph{Trouver la faille.} «~Notre routeur envoie chaque ordre passif vers la plateforme qui verse la plus forte rémunération, ce qui réduit les coûts de nos clients.~»
\end{exercise}

\section{Problème~: La rémunération qui a acheté une file}

\begin{problem}[{Problème du week-end --- la rémunération qui a acheté une file}]\label{pb:mx:fees-rebates-and-venue-economics:1}
Un teneur de marché cote une action sur une plateforme maker-taker et sur une plateforme inversée. Décider où afficher ses ordres et déterminer ce que finance la rémunération.

\medskip\noindent\textbf{Partie I --- \omterm{def:mx:fees-rebates-and-venue-economics:adjusted}{Prix nets des frais}.}
\begin{enumerate}
\item Donner les taux de création et de prise de la grille tarifaire de chaque plateforme.
\item Calculer le \omterm{def:mx:fees-rebates-and-venue-economics:adjusted}{prix net des frais} d'un achat de 1\,000 actions à 20.00 sur chacune.
\item Définir le \omterm{def:mx:fees-rebates-and-venue-economics:efftick}{pas de cotation effectif} et le calculer pour ces frais.
\item Pourquoi un \omterm{def:mx:the-limit-order-book:tick}{pas de cotation} discret pousse-t-il les plateformes à se livrer concurrence sur les frais~?
\end{enumerate}

\medskip\noindent\textbf{Partie II --- Neutralité.}
\begin{enumerate}[resume]
\item Énoncer et démontrer la proposition de neutralité.
\item Reprendre le transfert de 0.20 cent des frais de prise vers le côté création, avec et sans tick.
\item Où va la rente lorsque les cotations ne peuvent pas bouger~?
\item Quels constats Malinova et Park ont-ils faits sur la Toronto Stock Exchange~?
\end{enumerate}

\medskip\noindent\textbf{Partie III --- Deux plateformes.}
\begin{enumerate}[resume]
\item Décrire le marché simulé et les deux types de preneurs de liquidité.
\item Tous les preneurs étant sensibles aux frais, donner la part de volume de chaque plateforme, sa part exécutée et le temps médian jusqu'à la première exécution.
\item Donner la sélection adverse par action exécutée sur chaque plateforme et expliquer l'écart.
\item Donner la valeur par action affichée sur chaque plateforme.
\item Quelle rémunération rendrait un fournisseur indifférent, et MT peut-elle la verser~?
\end{enumerate}

\medskip\noindent\textbf{Partie IV --- Routage et verdict.}
\begin{enumerate}[resume]
\item Qu'est-ce qui change lorsque la moitié des preneurs dirigent leurs ordres selon la taille affichée~?
\item Et lorsqu'aucun d'entre eux n'est sensible aux frais~?
\item Qu'était le Transaction Fee Pilot, et que lui est-il arrivé~?
\item Quels constats Battalio, Corwin et Jennings ont-ils faits~?
\item Comment le nouveau plafond des frais d'accès modifie-t-il la marge de manœuvre de la plateforme maker-taker pour les rémunérations~?
\item Énoncer le \emph{résultat nommé}~: le taux d'exécution et le mark-out des ordres passifs sur la plateforme maker-taker et sur la plateforme inversée, ainsi que la rémunération à laquelle un
fournisseur de liquidité est indifférent.
\item En une phrase~: qu'achète une rémunération~?
\end{enumerate}
\end{problem}

\section{Questions d'entretien}

\begin{interviewq}[$\star$ \iqroles{trader}]\label{iq:mx:fees-rebates-and-venue-economics:1}
Quelle est la différence entre une plateforme maker-taker et une plateforme inversée, et où enverriez-vous un ordre passif que vous devez exécuter rapidement~?
\end{interviewq}

\begin{interviewq}[$\star\star$ \iqroles{trader, chercheur}]\label{iq:mx:fees-rebates-and-venue-economics:2}
La répartition des frais d'une bourse entre fournisseurs de liquidité et preneurs de liquidité est-elle \omterm{def:mx:fees-rebates-and-venue-economics:neutral}{neutre}~? Quand ne l'est-elle pas~?
\end{interviewq}

\begin{interviewq}[$\star\star$ \iqroles{chercheur}]\label{iq:mx:fees-rebates-and-venue-economics:3}
Les exécutions passives sur une plateforme ont des mark-out moins bons que sur une autre au même prix. Proposez deux explications et un test qui permette de les départager.
\end{interviewq}

\begin{interviewq}[$\star\star$ \iqroles{trader}]\label{iq:mx:fees-rebates-and-venue-economics:4}
Votre routeur classe les plateformes selon le \omterm{def:mx:fees-rebates-and-venue-economics:adjusted}{prix net des frais}. Quel effet cela a-t-il sur les teneurs de marché qui affichent leurs ordres sur la plateforme la plus chère à prendre~?
\end{interviewq}

\begin{interviewq}[$\star\star$ \iqroles{chercheur}]\label{iq:mx:fees-rebates-and-venue-economics:5}
Comment concevriez-vous une expérience pour mesurer l'effet des rémunérations sur la qualité de marché~?
\end{interviewq}

\begin{interviewq}[$\star\star\star$ \iqroles{trader}]\label{iq:mx:fees-rebates-and-venue-economics:6}
Le plafond des frais d'accès passe de 0.30 à 0.10 cent. Qu'advient-il des rémunérations maker-taker, des files d'attente et de votre P\&L en market making~?
\end{interviewq}
